\documentclass[12pt, a4paper]{article}
\RequirePackage[utf8]{inputenc}
\RequirePackage{multicol}
\RequirePackage{ragged2e}
\RequirePackage{graphicx}
\RequirePackage{amsmath}
\RequirePackage{xcolor}
\RequirePackage{geometry}
\RequirePackage{caption}
\RequirePackage{subcaption}
\RequirePackage{enumitem}
\RequirePackage{tikz}
\geometry{top = 2.5cm, bottom = 2.5cm, left = 1.5cm, right = 1.5cm}
\makeatletter
  \def\vhrulefill#1{\leavevmode\leaders\hrule\@height#1\hfill \kern\z@}
\makeatother

%-----------------------------------------------------------
%       TO START EDITING THE DOCUMENT JUMP TO LINE 148
%-----------------------------------------------------------
\title{Question Paper}
\author{}
\date{}
%------------------------------------------------------------
%                    HEADER INFORMATION
%------------------------------------------------------------
\newcommand{\hdr}[6]{
\begin{center}
    \ifx&#1& \else \leavevmode{\Huge #1}\\[5mm] \fi
    #2\\\vspace{5mm}
    {\Large \textbf{#3}}\\[2.5mm]
    {\Large \textbf{#4}}\\
\end{center}
\begin{raggedleft}
{\large \textbf{#5 \hfill #6}\\}\vspace{3mm} 
\vhrulefill{5pt}
\end{raggedleft}}

%------------------------------------------------------------
%                  GENERAL INSTRUCTIONS
%------------------------------------------------------------
\newcommand{\geninstructions}[1]{
\begin{raggedright}
\vspace{3mm}
\textbf{\large INSTRUCTIONS}
\textit{
\begin{itemize}
    #1
\end{itemize}}
\hrulefill\\
\hrulefill \\\vspace{5mm}
\end{raggedright}}

%-----------------------------------------------------------
%                   QUESTION SECTION
%-----------------------------------------------------------
\newcommand{\questionsection}[3]{
\vspace{2.5mm}
\begin{center}
    \textbf{\large \underline{SECTION - #1}}\\
{\small \textbf{#2}}
\end{center}
{\small 
\textit{\begin{itemize}[noitemsep, topsep=0pt]
    #3
\end{itemize}}}}

%-----------------------------------------------------------
%              4 OPTION MCQ SEPARATED IN 2 COLS
%-----------------------------------------------------------
\newcommand{\mcqfourtwo}[9]{
\vspace{2.5mm}
\begin{raggedright}
\textbf{Question #1:} #2 \hfill \textit{#7}\\
\textit{#8}\\
#9
\begin{multicols}{2}{}
(a) #3\\
(c) #5\\
\columnbreak
(b) #4\\
(d) #6\\
\end{multicols}
\end{raggedright}}

%-----------------------------------------------------------
%                   DESCRIPTIVE QUESTION
%-----------------------------------------------------------
\newcommand{\question}[4]{
\vspace{1.5mm}
\begin{raggedright}
\textbf{Question #1:} #2\hfill\textit{#3}\\
\textit{#4}\\
\end{raggedright}}

%-----------------------------------------------------------
%                 BEGIN DOCUMENT
%-----------------------------------------------------------

\begin{document}

\hdr{}{Class XI (ISC)}{Assignment - 1}{(Errors, Sig. Figs, Vectors \& Projectile Motion)}{1 Hr 30 Mins}{50 marks}

\geninstructions{
      \item All questions are compulsory.
      \item The paper consists of four sections: A, B, C, and D.
      \item Section A contains 10 MCQs of 1 mark each. Section B contains 5 questions of 2 marks each. Section C contains 5 questions of 3 marks each. Section D contains 3 long questions of 5 marks each.
      \item Use of a scientific calculator is permitted.
}

%-----------------------------------------------------------
%                     SECTION A
%-----------------------------------------------------------
\questionsection{A}{Answer all of the following multiple choice questions.}{
\item 10$\times$1=10 Marks
}

\mcqfourtwo{1}{The number of significant figures in the measurement $0.005020 \text{ kg}$ is:}
        {3}
        {4}
        {6}
        {7}
        {(1)}{}{}

\mcqfourtwo{2}{The dimensional formula for the Universal Gravitational Constant ($G$) is:}
        {$[M^{-1}L^{3}T^{-2}]$}
        {$[M^{1}L^{2}T^{-2}]$}
        {$[M^{-1}L^{2}T^{-2}]$}
        {$[M^{1}L^{3}T^{-2}]$}
        {(1)}{}{}

\mcqfourtwo{3}{A physical quantity $Z$ is given by $Z = \frac{A^2 B}{C^3}$. If the percentage errors in the measurement of $A, B$, and $C$ are $1\%, 2\%$, and $3\%$ respectively, the maximum percentage error in $Z$ is:}
        {$6\%$}
        {$13\%$}
        {$14\%$}
        {$8\%$}
        {(1)}{}{}

\mcqfourtwo{4}{If $\vec{A} \cdot \vec{B} = 0$ and both $\vec{A}$ and $\vec{B}$ are non-zero vectors, it implies that the vectors are:}
        {Parallel}
        {Anti-parallel}
        {Perpendicular}
        {Equal in magnitude}
        {(1)}{}{}

\mcqfourtwo{5}{For a given initial velocity, the angle of projection for which a projectile covers the maximum horizontal range is:}
        {$30^{\circ}$}
        {$45^{\circ}$}
        {$60^{\circ}$}
        {$90^{\circ}$}
        {(1)}{}{}

\mcqfourtwo{6}{Which of the following physical quantities has the dimensional formula $[M^1 L^{-1} T^{-2}]$?}
        {Force}
        {Work}
        {Pressure}
        {Power}
        {(1)}{}{}

\mcqfourtwo{7}{The magnitude of the cross product of two vectors is $\sqrt{3}$ times the magnitude of their dot product. The angle between the two vectors is:}
        {$30^{\circ}$}
        {$45^{\circ}$}
        {$60^{\circ}$}
        {$90^{\circ}$}
        {(1)}{}{}

\mcqfourtwo{8}{The minimum number of unequal coplanar vectors whose resultant can be zero is:}
        {Two}
        {Three}
        {Four}
        {Infinite}
        {(1)}{}{}

\mcqfourtwo{9}{A projectile is fired at an angle $\theta$ such that its maximum height is equal to its horizontal range ($H = R$). The angle of projection $\theta$ is:}
        {$\tan^{-1}(1)$}
        {$\tan^{-1}(2)$}
        {$\tan^{-1}(4)$}
        {$\tan^{-1}(1/4)$}
        {(1)}{}{}

\mcqfourtwo{10}{If $|\vec{A} \times \vec{B}|^2 + |\vec{A} \cdot \vec{B}|^2 = 144$ and $|\vec{A}| = 4$, then the magnitude of $\vec{B}$ is:}
        {3}
        {4}
        {12}
        {36}
        {(1)}{}{}

\newpage

%-----------------------------------------------------------
%                     SECTION B
%-----------------------------------------------------------
\questionsection{B}{Answer the following short answer questions.}{
\item 5$\times$2=10 Marks
}

\question{1}{State the rules for rounding off measurements ending in 5. Apply them to round off the values $3.785$ and $3.775$ to three significant figures.}
{(2)}{}

\question{2}{The mass of a solid sphere is measured as $(5.00 \pm 0.05) \text{ kg}$ and its radius as $(10.0 \pm 0.1) \text{ cm}$. Calculate the maximum percentage error in the determination of its density.}
{(2)}{}

\question{3}{Find a unit vector parallel to the resultant of the vectors $\vec{A} = 2\hat{i} + 3\hat{j} - \hat{k}$ and $\vec{B} = \hat{i} - 2\hat{j} + 3\hat{k}$.}
{(2)}{}

\question{4}{A projectile is fired at an angle of $30^{\circ}$ to the horizontal. Calculate the ratio of its maximum height attained to its horizontal range.}
{(2)}{}

\question{5}{Using the principle of homogeneity of dimensions, check the correctness of the equation $v = \sqrt{\frac{T}{m}}$, where $v$ is the velocity of a transverse wave on a string, $T$ is the tension in the string, and $m$ is the mass per unit length (linear mass density) of the string.}
{(2)}{}


%-----------------------------------------------------------
%                     SECTION C
%-----------------------------------------------------------
\questionsection{C}{Answer the following medium answer questions.}{
\item 5$\times$3=15 Marks
}

\question{1}{The period of oscillation of a simple pendulum is given by $T = 2\pi\sqrt{\frac{L}{g}}$. The measured value of $L$ is $20.0 \text{ cm}$ known to $1 \text{ mm}$ accuracy and the time for 100 oscillations of the pendulum is found to be $90 \text{ s}$ using a wrist watch of $1 \text{ s}$ resolution. Find the maximum percentage error in the determination of $g$.}
{(3)}{}

\question{2}{State the Parallelogram Law of Vector Addition. Derive the analytical expression for the magnitude of the resultant of two vectors $\vec{P}$ and $\vec{Q}$ inclined at an angle $\theta$ to each other.}
{(3)}{}

\question{3}{An FPV drone is flying horizontally at a steady speed of $20 \text{ m/s}$ at an altitude of $80 \text{ m}$. It releases a small environmental sensor payload to test a micro-ecosystem below. Neglecting air resistance, calculate: 
\newline (a) The time taken by the sensor payload to reach the ground.
\newline (b) The horizontal distance (range) from the drop point to where the sensor lands. (Take $g = 10 \text{ m/s}^2$).}
{(3)}{}

\question{4}{Given two vectors $\vec{A} = \hat{i} + 2\hat{j} + x\hat{k}$ and $\vec{B} = 2\hat{i} - \hat{j} + \hat{k}$. 
\newline (a) Find the value of $x$ for which the two vectors are exactly perpendicular to each other.
\newline (b) Using the calculated value of $x$, find the cross product $\vec{A} \times \vec{B}$.}
{(3)}{}

\question{5}{A particle is moving in a circular path. Assuming that the centripetal force ($F$) acting on it depends upon the mass ($m$) of the particle, its velocity ($v$), and the radius ($r$) of the circular path, derive an expression for the force using the method of dimensional analysis. (Take the dimensionless constant $k = 1$).}
{(3)}{}

\newpage

%-----------------------------------------------------------
%                     SECTION D
%-----------------------------------------------------------
\questionsection{D}{Answer the following long answer questions.}{
\item 3$\times$5=15 Marks
}

\question{1}{(i) For a projectile fired from the ground with an initial velocity $u$ at an angle $\theta$ with the horizontal, show mathematically that its trajectory is a parabola.
\newline (ii) Show that for a given initial velocity, there are two angles of projection $\theta_1$ and $\theta_2$ for which the horizontal range is the same. If $t_1$ and $t_2$ are the respective times of flight for these two angles, prove that $t_1 t_2 = \frac{2R}{g}$, where $R$ is the horizontal range.}
{(2+3)}{}

\question{2}{(i) A river $1 \text{ km}$ wide is flowing at a steady rate of $3 \text{ km/h}$. A swimmer, who can swim at a speed of $5 \text{ km/h}$ in still water, wishes to cross the river straight to the point exactly opposite his starting point. At what angle to the river flow should he swim?
\newline (ii) How much time will the swimmer take to cross the river along this shortest path?
\newline (iii) If instead, he dives strictly perpendicular to the river flow, how far downstream will he be carried by the time he reaches the opposite bank?}
{(2+1+2)}{}

\question{3}{(i) The Vander Waals equation for a real gas is given by:
$$ \left(P + \frac{a}{V^2}\right)(V - b) = RT $$
Where $P$ is pressure, $V$ is volume, $T$ is temperature, and $R$ is the universal gas constant. Deduce the dimensional formulas for the constants $a$ and $b$.
\newline (ii) The sum and difference of two non-zero vectors $\vec{A}$ and $\vec{B}$ are equal in magnitude, i.e., $|\vec{A} + \vec{B}| = |\vec{A} - \vec{B}|$. Prove mathematically that the vectors $\vec{A}$ and $\vec{B}$ must be strictly perpendicular to each other.}
{(3+2)}{}

\end{document}